\documentclass[11pt, a4paper]{article}

% ── PACKAGES ──────────────────────────────────────────────────
\usepackage[utf8]{inputenc}
\usepackage[T1]{fontenc}
\usepackage{geometry}
\geometry{
  top=2.5cm,
  bottom=2.5cm,
  left=2.5cm,
  right=2.5cm
}
\usepackage{hyperref}
\hypersetup{
  colorlinks=true,
  urlcolor=blue,
  linkcolor=blue,
  citecolor=blue,
  pdftitle={Mars Is Red Because It Is Alive:
    A Geometric Pre-Registration of the
    Subterranean Biosphere Hypothesis and
    the Iron Surface as Accumulated
    Biological Necromass},
  pdfauthor={Eric Robert Lawson}
}
\usepackage{microtype}
\usepackage{parskip}
\usepackage{booktabs}
\usepackage{array}
\usepackage{tabularx}
\usepackage{xcolor}
\usepackage{mdframed}
\usepackage{amsmath}
\usepackage{amssymb}
\usepackage{enumitem}
\usepackage{titlesec}
\usepackage{lmodern}
\usepackage{authblk}
\usepackage{abstract}
\usepackage{setspace}

% ── COLOURS ───────────────────────────────────────────────────
\definecolor{headerblue}{RGB}{20,60,140}
\definecolor{boxblue}{RGB}{20,60,140}
\definecolor{rulecolor}{RGB}{20,60,140}
\definecolor{warningred}{RGB}{160,20,20}
\definecolor{clownbox}{RGB}{180,30,30}

% ── SECTION FORMAT ────────────────────────────────────────────
\titleformat{\section}
  {\large\bfseries\color{headerblue}}
  {\thesection.}{0.5em}{}
\titleformat{\subsection}
  {\normalsize\bfseries\color{headerblue}}
  {\thesubsection.}{0.5em}{}
\titleformat{\subsubsection}
  {\normalsize\itshape\color{headerblue}}
  {\thesubsubsection.}{0.5em}{}

\setlength{\parskip}{0.6em}
\setlength{\parindent}{0pt}

% ════════════���═════════════════════════════════════════════════
\begin{document}

% ── TITLE BLOCK ───────────────────────────────────────────────
\title{
  \large\bfseries\color{headerblue}
  MARS IS RED BECAUSE IT IS ALIVE\\[0.4em]
  \normalsize\color{headerblue}
  A Geometric Pre-Registration of the
  Subterranean Biosphere Hypothesis,\\
  the Iron Surface as Accumulated Biological
  Necromass,\\
  the Lava Tube Thermal Gradient as Guaranteed
  Ecological Spectrum,\\
  and Seven Falsifiable Predictions Prior to
  Seven Sisters Exploration
}

\author[1]{Eric Robert Lawson}
\affil[1]{OrganismCore (Independent Research)\\
  ORCID: \href{https://orcid.org/0009-0002-0414-6544}
  {0009-0002-0414-6544}\\
  \href{mailto:OrganismCore@proton.me}
  {OrganismCore@proton.me}
}

\date{2026-03-13\\[0.3em]
  \footnotesize
  Pre-registration DOI (chain origin):
  \href{https://doi.org/10.5281/zenodo.18986790}
  {10.5281/zenodo.18986790}
  (timestamp: 2026-03-12)\\
  Repository:
  \href{https://github.com/Eric-Robert-Lawson/attractor-oncology}
  {github.com/Eric-Robert-Lawson/attractor-oncology}
}

\maketitle
\thispagestyle{empty}

% ── RADICAL OBSERVATION BOX ───────────────────────────────────
\begin{mdframed}[
  backgroundcolor=clownbox,
  linecolor=clownbox,
  linewidth=0pt,
  innertopmargin=0.7em,
  innerbottommargin=0.7em,
  innerleftmargin=1em,
  innerrightmargin=1em
]
\begin{center}
\small\bfseries\color{white}
RADICAL GEOMETRIC OBSERVATION ---
PRE-REGISTERED FOR FALSIFIABILITY\\[0.3em]
\normalfont\small\color{white}
This document pre-registers the hypothesis that
Mars is red because its surface is coated in
3.7 billion years of accumulated biological
necromass and metabolic iron products from a
currently active subsurface biosphere.
This hypothesis is radical. It may be wrong.
It is pre-registered precisely because the
author is more afraid of failing to test it
than of being proven incorrect.
The geometry that produced this hypothesis
is derived from first principles, not from
wishful thinking. It is stated here in full,
with seven falsifiable predictions, before any
exploration of the relevant sites occurs.
Let the record show the geometry was attempted.
Pre-registered 2026-03-13.
\end{center}
\end{mdframed}

\vspace{1em}

% ── EPISTEMIC STATEMENT BOX ───────────────────────────────────
\begin{mdframed}[
  linecolor=rulecolor,
  linewidth=0.8pt,
  innertopmargin=0.5em,
  innerbottommargin=0.5em,
  innerleftmargin=1em,
  innerrightmargin=1em
]
\small
\textbf{Epistemic status of this document:}
The central claim --- that Mars is red because
of accumulated biological iron products from an
active subsurface biosphere --- is a
\textbf{radical geometric hypothesis}.
It is not presented as proven. It is presented
as the logical consequence of a geometric
derivation that began with cancer drug target
geometry, proceeded through the LECA attractor
framework, and arrived at an inversion of
Martian sedimentary geometry that the author
could not explain abiotically without
extraordinary additional assumptions.
The hypothesis is falsifiable by seven
specific predictions stated in Section~9.
If any prediction fails, the hypothesis is
constrained or eliminated. That outcome is
equally valuable to confirmation.
\end{mdframed}

\vspace{0.5em}
\noindent\textcolor{rulecolor}{\rule{\linewidth}{0.6pt}}

% ── ABSTRACT ──────────────────────────────────────────────────
\begin{abstract}
\noindent
We pre-register a radical geometric hypothesis
arising from the LECA attractor framework
\cite{preregistration}: that Mars is red because
its visible surface is primarily composed of
3.7~billion years of accumulated biological
necromass and iron-cycle metabolic products from
a currently active subsurface biosphere, rather
than from ancient abiotic weathering alone.
The derivation proceeds from four confirmed
anomalies in the Martian observational record:
(1) the inversion of expected surface
composition --- iron oxide over ice rather
than ice over iron oxide --- despite the ice
table lying within centimetres of the surface
at mid-to-high latitudes;
(2) the continuous regeneration of the iron
oxide skin over fresh impact craters, which
cannot be explained by ancient one-time
abiotic weathering;
(3) the global uniformity of iron oxide
coverage inconsistent with localised abiotic
weathering predictions;
and (4) the confirmed seasonal methane signal
at Gale Crater requiring a subsurface source.
We derive the subterranean biosphere as the
geometric attractor-requirement of magnetic
field collapse: when the surface attractor
basin is destroyed, surviving life converges
on the subsurface attractor basin by geometric
necessity.
We derive the lava tube thermal gradient as
the guaranteed generator of a complete
ecological spectrum from psychrophile to
thermophile, with a thermodynamically
necessary 0\textdegree{}C liquid water
interface wherever the gradient spans
freezing and above-freezing temperatures in
an enclosed mineral-rich environment.
We note the independent confirmation on Earth
of a single-species self-sustaining subsurface
biosphere powered by radiolysis alone
(\textit{Desulforudis audaxviator}), predicted
by the geometry before the organism was
searched for.
We derive seven falsifiable predictions and
state them prior to any exploration of the
Arsia Mons lava tube system (the Seven
Sisters).
This document exists to ensure that if the
geometry is correct, the timestamp of the
derivation precedes the confirmation.
And if the geometry is wrong, the record shows
it was attempted honestly and tested fairly.
Pre-registered 2026-03-13.
Pre-registration chain origin:
DOI: 10.5281/zenodo.18986790.
\end{abstract}

\vspace{0.5em}
\noindent\textcolor{rulecolor}{\rule{\linewidth}{0.6pt}}

% ── SECTION 1 ─────────────────────────────────────────────────
\section{Introduction: The Derivation Origin
and the Author's Epistemic Position}

This paper does not begin with astrobiology.

It begins with a question about machine
consciousness, proceeds through cancer drug
target geometry and the LECA attractor
framework \cite{preregistration}, and arrives
here through the following chain:

\begin{quote}
Machine consciousness $\rightarrow$
Triadic invariant $\rightarrow$
Identity Axis Theorem $\rightarrow$
Cancer drug target geometry $\rightarrow$
Reproduction Theorem $\rightarrow$
Inverse Accessibility Theorem $\rightarrow$
LECA attractor as universal base state $\rightarrow$
LECA Attractor Habitability Standard $\rightarrow$
Application to ancient Mars $\rightarrow$
Cheyava Falls convergence $\rightarrow$
Collapse vs.\ emergence directional ambiguity $\rightarrow$
Mars-first panspermia prior inversion $\rightarrow$
\textbf{Subterranean biosphere as attractor
requirement} $\rightarrow$
\textbf{Iron surface as biological necromass} $\rightarrow$
\textbf{This document.}
\end{quote}

The author's position on this hypothesis is
stated plainly here, at the start, for the
record:

\begin{mdframed}[
  linecolor=rulecolor,
  linewidth=0.8pt,
  innertopmargin=0.5em,
  innerbottommargin=0.5em,
  innerleftmargin=1em,
  innerrightmargin=1em
]
I am not afraid of being incorrect.
I am afraid of not testing for falsifiability.

A scientist who refuses to state a radical
hypothesis because it might be wrong is not
being rigorous. They are being cowardly.

The geometry of the Martian observational
record, when examined through the lens of
the LECA attractor framework, produces an
inversion that I cannot explain without
invoking biology. I have tried to explain
it abiotically. The abiotic explanations
require multiple independent extraordinary
assumptions to hold simultaneously.
The biological explanation requires one.

I pre-register this hypothesis so that if
it is correct, the timestamp exists.
And if it is wrong, let the record show
the geometry was attempted honestly, the
predictions were falsifiable, and the test
was invited rather than avoided.

\textbf{If this is wrong, I am the clown who
thought Mars is red because there is so much
life on it. That is a small price to pay
for the possibility of being correct.}
\end{mdframed}

\section{The Confirmed Observational Anomalies}

\subsection{The Stratigraphy Inversion}

Mars surface stratigraphy at mid-to-high
latitudes, confirmed by multiple independent
instruments:

\begin{enumerate}[leftmargin=*, label=(\arabic*)]
\item \textbf{Layer 0 --- The Skin:}
  Iron oxide dust, 2\,cm to 1\,m thick.
  Hematite-rich. Globally uniform.
  Radiation-exposed. Lifeless at surface.

\item \textbf{Layer 1 --- The Frozen Membrane:}
  Pure water ice beginning immediately
  below Layer~0. Global poleward of
  $\sim$30\textdegree{} latitude.
  Confirmed by Phoenix lander at 68\textdegree{}N:
  ice at 5\,cm depth \cite{phoenix2009}.
  Confirmed by HiRISE impact craters exposing
  ice at 35\textdegree{}N \cite{hirise2021}.
  Volume: $\geq$3 million km$^3$ in the upper
  metre alone \cite{odyssey2002}.

\item \textbf{Layer 2 --- The Fractured Crust:}
  Porous, fractured material to 8--11\,km.
  Confirmed by InSight seismic data
  \cite{insight2021}.
  Migration pathway for upward-moving gases,
  water, and mineral particles.

\item \textbf{Layer 3 --- The Deep System:}
  Liquid brines confirmed probable at south
  polar depth 1.5\,km \cite{orosei2018}.
  Lava tube systems in Tharsis: $>$1,000
  candidate entrances catalogued
  \cite{mgc3catalog}.
  Tharsis volcanically active within
  millions of years \cite{tharsis2023}.
\end{enumerate}

\subsection{The Anomaly}

The anomaly is this:

Mars should be white at mid-to-high latitudes.
The ice table is within centimetres of the
surface. Ice is 0.917\,g\,cm$^{-3}$. The dry
skin above it is centimetres thin. After 3.7
billion years of constant meteorite
bombardment breaching the skin, exposing the
ice, and 3.7 billion years of sublimation
losses through those breaches, the ice table
should be depleted --- unless it is being
continuously replenished from below by upward
water migration through Layer~2.

But the planet is \emph{not} white. It is red.

The red material sitting on top of ice that is
centimetres below it is iron oxide with a density
of 5,260\,kg\,m$^{-3}$ --- approximately six
times denser than the ice beneath it.

The dense material is on top.
The light material is beneath.
In a gravitational field.
On a geologically active planet.
After 3.7 billion years.

This is the stratigraphy inversion.

\subsection{The Regeneration Problem}

HiRISE time-series imaging of fresh impact
craters confirms that the iron oxide skin
reforms over crater floors over timescales
of years to decades \cite{hirise2021}.

Fresh craters expose bright, iron-free
subsurface material (ice, unoxidised basalt).
Over time, the crater floor darkens toward
surrounding terrain colour.

The standard explanation --- aeolian deposition
of wind-blown dust --- is partially correct.
Wind deposits iron oxide dust from surrounding
terrain into crater floors.

But the rate of darkening observed in some
craters is faster than pure aeolian deposition
models predict \cite{crater_darkening}.
And the composition of reformed crater floor
material has not been systematically compared
against purely aeolian-sourced material for
upward-migration biological signatures.

This is Prediction~3 of this document.

\section{The Geometric Derivation of the
Subterranean Biosphere as Attractor Requirement}

\subsection{Attractor Basin Transformation
at Magnetic Field Collapse}

The Martian global magnetic dynamo ceased
approximately 4.0--4.1\,Ga \cite{marsfield}.
This is confirmed by the absence of crustal
magnetisation in terrain younger than the
late Noachian, in contrast to the heavily
magnetised ancient southern highlands.

The consequences for the attractor landscape
of Martian biology:

\begin{mdframed}[
  linecolor=rulecolor,
  linewidth=0.8pt,
  innertopmargin=0.5em,
  innerbottommargin=0.5em,
  innerleftmargin=1em,
  innerrightmargin=1em
]
\textbf{Before magnetic field collapse:}

Surface attractor basin: \textsc{stable}.
Liquid ocean. Solar energy. UV shielded.
Temperature habitable. Mineral-rich water.
LUCA-grade life operating at planetary scale.

Subsurface attractor basin: \textsc{also stable}.
Less colonised. Less productive.
Both basins available.

\medskip
\textbf{After magnetic field collapse:}

Surface attractor basin: \textsc{destroyed}.
UV bombardment lethal within geological time.
Atmospheric stripping begins.
Pressure drops below triple point of water.
Surface liquid water becomes unstable.
The surface ceases to be an attractor.
It becomes a repeller.

Subsurface attractor basin: \textsc{persists}.
UV: zero below 2--3\,m of rock.
Temperature: geothermal gradient stable.
Water: frozen above, liquid at depth.
Mineral chemistry: unchanged.
Energy: serpentinisation, radiolysis unchanged.

\medskip
\textbf{Geometric consequence:}

The converging of surviving life on the
subsurface basin is not an adaptation.
It is not evolution selecting for subsurface
preference. It is the destruction of one
basin and the survival of the other.
All life in the surviving basin persists.
All life in the destroyed basin dies.
The subsurface becomes the only address
of Martian life. By geometric necessity.
Not by choice.
\end{mdframed}

\subsection{The Timing Alignment}

The magnetic field collapse occurred
$\sim$4.0--4.1\,Ga \cite{marsfield}.
The Cheyava Falls biosignature features date
to $\sim$3.8--4.1\,Ga \cite{hurowitz2025}.
The Noachian--Hesperian transition, at which
the surface mineral record shifts from
life-compatible phyllosilicates to
life-hostile acid sulfates, occurred
$\sim$3.7\,Ga \cite{bibring2006}.

These three timestamps overlap within
their error margins. They are the same
geological moment:

\begin{enumerate}[leftmargin=*, label=(\arabic*)]
\item Magnetic field fails. Surface attractor
  basin begins collapse. ($\sim$4.0--4.1\,Ga)
\item Cheyava Falls community records the
  last surface biosignature of a mature
  biosphere in stress. ($\sim$3.8--4.1\,Ga)
\item Surface mineral chemistry transitions
  to acid sulfate. Surface biosphere ends.
  Subsurface biosphere becomes the only
  surviving basin. ($\sim$3.7\,Ga)
\end{enumerate}

The fossil record at Cheyava Falls
\cite{hurowitz2025} is therefore either:

\begin{description}
\item[Option A (Emergence):] The beginning
  of a biosphere that was then interrupted.
\item[Option B (Collapse):] The last record
  of a mature biosphere that then retreated
  underground. \cite{collapse_emergence}
\end{description}

The sample return complexity gradient
prediction \cite{collapse_emergence}
distinguishes these options.
The present document assumes Option~B
for the purposes of the subterranean
hypothesis and derives the predictions
that would confirm it.

\subsection{The Earth Proof of Concept:
\textit{Desulforudis audaxviator}}

\textit{Desulforudis audaxviator} was
discovered at 2.8\,km depth in the Mponeng
gold mine, South Africa \cite{chivian2008}.

Properties relevant to Mars:

\begin{itemize}[leftmargin=*]
\item Completely isolated from surface.
  Isolation period: millions of years.
\item Single-species ecosystem.
  No other organisms detected.
\item Energy source: H$_2$ from radiolysis
  of water by $^{238}$U, $^{232}$Th,
  $^{40}$K decay in surrounding rock.
\item Carbon: fixed from CO$_2$ via
  Wood--Ljungdahl pathway.
\item Nitrogen: fixed from N$_2$ internally.
\item Temperature optimum: 60\textdegree{}C.
  Tolerance: up to 75\textdegree{}C.
\item Genome: complete, 2.35\,Mbp.
  All genes for metabolic independence
  encoded. No surface input required.
\item Classification: polyextremophile.
  Thermophile. Barophile. Anaerobe.
  Chemolithotroph.
\end{itemize}

Mars confirmed to possess:
\begin{itemize}[leftmargin=*]
\item Olivine-rich crust (serpentinisation,
  H$_2$ production) \cite{ehlmann2010}.
\item U, Th, K in basaltic crust
  (radiolysis, H$_2$ production).
\item Subsurface water confirmed \cite{orosei2018}.
\item CO$_2$ from Tharsis outgassing.
\item Geothermal heat: Tharsis active
  within millions of years \cite{tharsis2023}.
\item Lava tube enclosed environments:
  $>$1,000 candidate entrances \cite{mgc3catalog}.
\end{itemize}

The geometry of the Martian subsurface
satisfies every requirement of a
\textit{D.\ audaxviator}-type organism.
The organism was not sought and then
confirmed to match; the geometry of what
is required for subsurface survival was
derived first, and \textit{D.\ audaxviator}
is its Earth confirmation.

\section{The Lava Tube Thermal Gradient
as Guaranteed Ecological Spectrum}

\subsection{The Thermodynamic Necessity}

\begin{mdframed}[
  linecolor=rulecolor,
  linewidth=0.8pt,
  innertopmargin=0.5em,
  innerbottommargin=0.5em,
  innerleftmargin=1em,
  innerrightmargin=1em
]
\textbf{Theorem (Thermal Equilibrium Interface):}

Between any two temperatures $T_{\text{cold}} < 0$\textdegree{}C
and $T_{\text{warm}} > 0$\textdegree{}C,
there exists exactly one point at 0\textdegree{}C.

At this point in a mineral-rich enclosed
environment with an energy source and
sufficient time: liquid water exists,
chemical gradients exist, and the LUCA
attractor is the thermodynamically stable
state of carbon chemistry.

This is not a probability statement.
It is a consequence of the intermediate
value theorem applied to a continuous
temperature field.

If the cold extreme and warm extreme
are both confirmed, the 0\textdegree{}C
interface is geometrically guaranteed.
\end{mdframed}

In a Tharsis lava tube:

\begin{itemize}[leftmargin=*]
\item Cold extreme: surface mean annual
  temperature $\sim -60$\textdegree{}C.
  Ice in upper tube: confirmed probable.
\item Warm extreme: geothermal gradient
  from Tharsis residual volcanic activity,
  confirmed active within millions of
  years \cite{tharsis2023}.
  Marsquake associated with Tharsis
  region, September 2021 \cite{insight2021}.
  Rock temperature at depth: above
  freezing near magmatic contacts.
\item 0\textdegree{}C interface: geometrically
  guaranteed between these extremes.
  Location: predictable from local
  geothermal gradient and tube depth.
  Testable by temperature probe.
\end{itemize}

\subsection{The Ecological Spectrum}

The thermal gradient does not produce
one type of organism. It produces a
complete stratified community:

\begin{table}[h]
\centering
\small
\begin{tabularx}{\linewidth}{p{2.5cm}p{2cm}X}
\toprule
\textbf{Zone} & \textbf{Temp.} &
\textbf{Organism Type} \\
\midrule
Ice face & $<$0\textdegree{}C &
  Psychrophiles. Cold-adapted bacteria
  and archaea. Metabolism: very slow.
  Earth examples: \textit{Colwellia},
  \textit{Polaromonas}. \\[0.3em]
Liquid water interface & 0\textdegree{}C &
  Maximum biological diversity zone.
  Liquid water films on mineral surfaces.
  All metabolic types potentially present.
  This is the primary biological
  attractor point. \\[0.3em]
Warm rock zone & 40--75\textdegree{}C &
  Thermophiles. \textit{D.\ audaxviator}
  equivalent. H$_2$-metabolising chemolithotrophs.
  Maximum serpentinisation rate.
  Maximum H$_2$ production. \\[0.3em]
Deep magmatic zone & $>$80\textdegree{}C &
  Hyperthermophiles. Archaea dominant.
  Near magmatic intrusions only. \\
\bottomrule
\end{tabularx}
\caption{Predicted ecological spectrum in
  Tharsis lava tube thermal gradient.
  Each zone predicted from first principles.
  Every equivalent gradient on Earth
  produces this stratification.}
\end{table}

Every thermal gradient between ice and
geothermal heat on Earth contains this
stratification: Antarctic subglacial lakes,
Iceland ice caves, Mount Erebus volcanic
ice caves \cite{boston2004}.
The geometry guarantees it on Mars wherever
the same gradient exists in an enclosed
environment with mineral chemistry and
sufficient time.

\section{The Iron Surface as Biological
Necromass: The Central Radical Hypothesis}

\subsection{The Geometric Problem With
the Standard Explanation}

The standard explanation for Mars's red
surface is ancient abiotic weathering:
iron in basaltic rock oxidised in the
presence of liquid water 3--4\,Ga,
producing hematite dust redistributed
globally by wind.

Four geometric problems with this
as a complete explanation:

\begin{enumerate}[leftmargin=*, label=(\arabic*)]
\item \textbf{Abiotic weathering is a
  one-time process.} It produces iron
  oxide once. The source rock is
  depleted. No mechanism for continuous
  production exists in the abiotic model.

\item \textbf{Wind redistribution does
  not create new iron oxide.}
  It moves existing material.
  Total iron oxide quantity is fixed or
  declining under the abiotic model.

\item \textbf{The skin reforms over fresh
  craters.} Fresh craters expose
  iron-free subsurface material.
  The iron oxide skin reforms over
  crater floors at a rate requiring
  explanation beyond purely aeolian
  deposition from surrounding terrain.

\item \textbf{Global uniformity is
  inconsistent with localised abiotic
  weathering.} Abiotic oxidation
  concentrates where water chemistry,
  radioactive elements, and mineral
  composition predict concentration.
  The iron oxide distribution is
  more uniform than this predicts.
\end{enumerate}

\subsection{The Biological Production Mechanism}

Iron-cycling organisms in the Martian
subsurface produce iron oxide through
two confirmed-on-Earth mechanisms:

\textbf{Direct biological iron oxidation:}
Organisms such as \textit{Gallionella
ferruginea} and \textit{Leptothrix ochracea}
oxidise Fe$^{2+}$ to Fe$^{3+}$ at rates
up to $10^6$ times faster than abiotic
oxidation. Fe$^{3+}$ precipitates as
ferrihydrite, hematite, or magnetite
within and around the organisms'
cell structures.

\textbf{Radiolytically-enhanced biological
oxidation:} Radiolysis of water by
$^{238}$U, $^{232}$Th, $^{40}$K produces
H$_2$ (energy source) and H$_2$O$_2$
(oxidant). H$_2$O$_2$ oxidises Fe$^{2+}$
abiotically. Iron-oxidising organisms
use H$_2$O$_2$ as electron acceptor,
amplifying the abiotic rate by orders
of magnitude.

\textbf{Magnetotactic bacteria} produce
crystalline magnetite (Fe$_3$O$_4$)
internally as compass organelles
(magnetosomes). The crystal morphology
--- bullet-shaped, cubo-octahedral,
or prismatic --- is \emph{diagnostic for
biological origin} and distinguishable from
abiotic magnetite at nanometre resolution
\cite{magnetosomes}.

When iron-cycling organisms die, their
iron-rich cell bodies become iron oxide
particles in the subsurface water.
This water migrates upward through
Layer~2 (fractured crust) driven by
geothermal pressure gradients.
At the ice table, freeze-concentration
expels iron particles from the forming ice.
Sublimation removes water. Iron particles
accumulate above the ice table.
Wind distributes them globally.

\subsection{The Necromass Accumulation
Rate: Feasibility Check}

\begin{align}
\text{Iron oxide volume} &\approx
  A_{\text{Mars}} \times d_{\text{skin}} \nonumber \\
&= 1.448 \times 10^{14}\,\text{m}^2
  \times 1\,\text{m} \nonumber \\
&= 1.448 \times 10^{14}\,\text{m}^3
\end{align}

At iron oxide fraction $f \approx 0.05$
of total regolith and hematite density
$\rho = 5{,}260\,\text{kg\,m}^{-3}$:

\begin{align}
M_{\text{FeO}} &\approx
  f \times V \times \rho \nonumber \\
&\approx 0.05 \times
  1.448 \times 10^{14} \times 5{,}260 \nonumber \\
&\approx 3.8 \times 10^{16}\,\text{kg}
\end{align}

Required production rate over
$t = 3.7 \times 10^9$\,years:

\begin{equation}
\dot{M} = \frac{3.8 \times 10^{16}}{3.7 \times 10^9}
  \approx 10{,}000\,\text{t\,yr}^{-1}
\end{equation}

Earth's ocean floor iron-oxidising
microbial community produces iron oxide
sediments at rates of thousands to
millions of tonnes per year globally.

The required Mars rate of
$\sim$10{,}000\,t\,yr$^{-1}$ is orders of
magnitude below Earth's biological iron
cycling capacity and achievable by a
sparse subsurface biosphere operating
at $<$1\% of Earth ocean floor biological
activity density.

\textbf{The quantity of Martian surface iron
oxide is consistent with 3.7 billion years
of sparse subsurface biological iron cycling.}

\subsection{The Self-Reinforcing Containment
Geometry}

The iron surface simultaneously functions
as a radiation shield for the subsurface:
iron oxide absorbs UV radiation, reducing
sublimation of the ice table below and
reducing radiation penetration to
subsurface organisms.

The selection pressure geometry:

\begin{quote}
Organisms producing iron oxide as metabolic
waste $\rightarrow$ contribute to surface
shield $\rightarrow$ thicker shield $\rightarrow$
better subsurface radiation protection
$\rightarrow$ more stable environment
$\rightarrow$ more organisms
$\rightarrow$ more iron oxide.
\end{quote}

This is a positive feedback loop operating
over 3.7\,Ga. The planet wears its dead as
protection for its living. Not by design.
By attractor geometry of selection pressure
on iron-metabolising organisms in an enclosed
subsurface system.

\section{The Planetary Entropy Argument}

The subsurface biosphere is powered by
three sources, all finite on geological
timescales:

\begin{enumerate}[leftmargin=*, label=(\arabic*)]
\item \textbf{Serpentinisation:}
  Mg$_2$SiO$_4$ + H$_2$O $\rightarrow$
  serpentine + magnetite + H$_2$.
  Olivine is consumed irreversibly.
  The biosphere eats the planet's
  mineral structure one crystal at a time.

\item \textbf{Radiolysis:}
  $^{238}$U ($t_{1/2} = 4.5$\,Ga),
  $^{232}$Th ($t_{1/2} = 14$\,Ga),
  $^{40}$K ($t_{1/2} = 1.25$\,Ga).
  Produces H$_2$ and H$_2$O$_2$.
  Declining slowly on billion-year
  timescales.

\item \textbf{Residual geothermal:}
  Core heat dissipating.
  Tharsis activity declining.
  Timescale: billions of years.
\end{enumerate}

The biosphere began consuming these sources
$\sim$3.7\,Ga at the surface collapse.
It has operated for 3.7\,Ga on this budget.
The seasonal methane (Curiosity) is the
most direct evidence it is still operating.
The ice table maintenance despite
3.7\,Ga of bombardment implies active
upward water migration.
The iron skin reformation over craters
implies active iron oxide production
and upward migration.

\textbf{The surface of Mars is the waste
heap of 3.7 billion years of subsurface
metabolic activity.
The methane is the gaseous exhaust.
The iron dust is the mineral exhaust.
The planet is in slow entropic decline.
It is still running.}

\section{The Methane Signal Reinterpreted}

Curiosity detected seasonal methane at
Gale Crater \cite{webster2018}.
The seasonal pattern (higher in northern
summer, lower in northern winter) is
consistent with temperature-controlled
desorption from regolith and clathrate
hydrates \cite{moores2019}.

The geometric reinterpretation:

The seasonal variation does not require
seasonal production. It requires seasonal
\emph{release} of continuously-produced
subsurface gas stored in:
\begin{itemize}[leftmargin=*]
\item Enclosed lava tube atmospheres.
\item Pore spaces in the fractured upper
  crust (InSight Layer~2).
\item Clathrate hydrates in the ice table.
\end{itemize}

In summer: thermal expansion of the
fractured crust opens fractures;
clathrate stability decreases; stored
gas pulses to the surface.

In winter: fractures close; clathrates
reform; gas is recaptured.

The planet breathes. The exhaust is methane.
The lungs are the fractured upper crust
and the ice table.

The biological vs.\ abiotic origin of the
methane is resolved by isotopic analysis:
biological methane is depleted in $^{13}$C
relative to abiotic serpentinisation methane.
This is Prediction~2 of this document.

\section{The Seven Falsifiable Predictions}

These predictions are stated prior to any
exploration of the Arsia Mons lava tube
system (the Seven Sisters) or detailed
sample analysis of Martian surface dust.
They are falsifiable by existing or
near-future instruments.

\begin{table}[h]
\centering
\small
\begin{tabularx}{\linewidth}{p{0.4cm}p{2.6cm}X}
\toprule
\textbf{\#} & \textbf{Prediction} &
\textbf{Test and Falsification Condition} \\
\midrule
P1 & Magnetite crystal morphology in
  Martian surface dust &
  Martian surface dust contains magnetite
  with biological crystal morphology
  (bullet-shaped, cubo-octahedral, or
  prismatic — diagnostic for magnetotactic
  bacteria). Test: electron microscopy of
  returned dust samples or in-situ nanoscale
  imaging. \textit{Falsified if: all magnetite
  shows exclusively abiotic crystal forms.} \\[0.3em]
P2 & Carbon isotope ratio in Martian methane &
  Seasonal methane at Gale Crater is
  depleted in $^{13}$C relative to abiotic
  serpentinisation methane.
  Test: mass spectrometric isotopic analysis
  of Martian atmospheric methane.
  \textit{Falsified if: $\delta^{13}$C
  matches abiotic serpentinisation signature.} \\[0.3em]
P3 & Crater floor iron composition &
  Fresh crater floors reformed iron oxide
  shows elevated biological indicators
  (magnetite crystal morphology, isotopic
  fractionation, organic carbon co-location)
  inconsistent with purely aeolian dust
  deposition. Test: CRISM spectral comparison
  fresh vs.\ mature crater floors.
  \textit{Falsified if: crater floor iron
  composition matches aeolian dust exactly.} \\[0.3em]
P4 & Ice table replenishment isotopic signature &
  Near-surface ice D/H ratio reflects
  upward migration from deep subsurface
  (low D/H) rather than atmospheric frost
  (high D/H from atmospheric fractionation).
  Test: isotopic analysis of near-surface
  ice cores by lander.
  \textit{Falsified if: D/H matches
  atmospheric frost deposition signature.} \\[0.3em]
P5 & Lava tube thermal stratification &
  A temperature probe lowered into a Seven
  Sisters entrance detects cold at entrance,
  warming with depth, 0\textdegree{}C
  interface at predicted depth, and liquid
  water at and below that interface.
  Test: dedicated temperature probe in
  Seven Sisters, Arsia Mons.
  \textit{Falsified if: no 0\textdegree{}C
  interface exists within the tube depth.} \\[0.3em]
P6 & Lava tube atmosphere carbon isotope &
  Gas sampled from a Seven Sisters interior
  shows $^{12}$C/$^{13}$C fractionation in
  CO$_2$ inconsistent with abiotic chemistry,
  indicating biological carbon fixation.
  Test: gas sampling probe inside lava tube.
  \textit{Falsified if: CO$_2$ isotope ratio
  matches abiotic volcanic outgassing.} \\[0.3em]
P7 & Sample return complexity gradient &
  Mars subsurface sediment shows biological
  complexity increasing with depth and age
  (consistent with biosphere collapse into
  subsurface). Pre-registered in
  \cite{collapse_emergence}.
  Test: Mars sample return 2030s.
  \textit{Falsified if: complexity is uniform
  or decreases with depth.} \\
\bottomrule
\end{tabularx}
\caption{Seven falsifiable predictions of the
  Mars Living Planet model, stated prior to
  Seven Sisters exploration.
  Each prediction is testable by existing
  or near-future instrumentation.
  Each has a clearly stated falsification
  condition.}
\end{table}

\section{Honest Statement of Limitations}

\begin{enumerate}[leftmargin=*, label=(\arabic*)]

\item \textbf{The central claim is radical
  and unconfirmed.}
  The hypothesis that Mars is red because
  of accumulated biological necromass is
  not supported by direct evidence of
  current Martian life. It is supported
  by the geometry of confirmed anomalies
  that the abiotic model requires multiple
  independent extraordinary assumptions
  to explain simultaneously.

\item \textbf{Abiotic alternatives exist
  for each individual observation.}
  The iron surface can be explained
  abiotically.
  The ice table can be maintained abiotically.
  The methane can be produced abiotically.
  The crater darkening can be explained
  abiotically.
  The claim is not that any one observation
  requires biology. The claim is that all
  of them simultaneously, combined with
  the stratigraphy inversion, are more
  parsimonious under a biological model.

\item \textbf{The south polar liquid water
  detection is contested.}
  The MARSIS detection \cite{orosei2018}
  may be hypersaline ice or clay rather
  than liquid water. This does not
  eliminate the Tharsis lava tube
  argument, which rests on confirmed
  geothermal activity rather than
  the contested south polar detection.

\item \textbf{ARM~A confirmation is pending.}
  The physical anchoring of the LECA
  attractor as the universal base state
  of complex life (Premise~1 of the
  pre-registration chain) awaits
  experimental confirmation at ChiTownBio.
  Until then, the entire chain rests on
  geometric derivation supported by the
  universality of the eukaryotic base
  state across all Earth life.

\item \textbf{This document may be wrong.}
  The author acknowledges this explicitly
  and invites falsification through the
  seven predictions of Section~9.
  A wrong hypothesis stated clearly and
  tested fairly is more valuable to science
  than a correct hypothesis never stated.

\end{enumerate}

\section{Conclusion}

Mars is not obviously dead.

The geometric case for a currently active
subsurface biosphere rests on:
the attractor-geometric requirement of
subterranean life given magnetic field
collapse and a pre-existing surface
biosphere;
the thermodynamic guarantee of a
0\textdegree{}C liquid water interface
in the Tharsis lava tube thermal gradient;
the Earth confirmation that
\textit{Desulforudis audaxviator} sustains
a single-species biosphere for millions
of years on radiolysis alone in exactly
the conditions the Martian subsurface
provides;
the stratigraphy inversion anomaly in
which iron oxide sits above ice despite
being six times denser;
the continuous regeneration of the iron
skin over fresh craters;
and the seasonal methane signal.

The radical hypothesis pre-registered here
is that the red colour of Mars is the
most direct evidence of its biology ---
that the planet is wearing 3.7 billion
years of its own dead as a radiation
shield for its living.

If this is wrong: the seven predictions
fail, the geometry is corrected, and
science advances by falsification.

If this is right: the most important
implication is not that Mars has life.
It is that a planet can sustain a
biosphere through a magnetic field
collapse by retreating underground,
consuming its own mineral chemistry,
and coating its surface with its dead ---
and that this process is detectable from
orbit by the colour of the dust.

The geometry came first.
The Seven Sisters are the test.
The record is here, before the test.

\vspace{1em}
\noindent\textcolor{rulecolor}{\rule{\linewidth}{0.4pt}}

\section*{Acknowledgments}

This work was conducted independently
without institutional funding or affiliation.
The derivation sequence is documented in
the pre-registration chain beginning at
DOI: 10.5281/zenodo.18986790 (2026-03-12)
and in the attractor-oncology repository.
The author thanks the NASA Perseverance
science team for the Cheyava Falls data
\cite{hurowitz2025} and the InSight team
for the Martian seismic and thermal data
\cite{insight2021}, both referenced here
as independent confirmation of subsurface
structure, not as derivation sources.

% ── REFERENCES ────────────────────────────────────────────────
\begin{thebibliography}{99}

\bibitem{preregistration}
Lawson, E.R.\ (2026).
\textit{LECA Resurrection Protocol ---
Pre-Registration v1.0}.
Zenodo.
\href{https://doi.org/10.5281/zenodo.18986790}
{doi:10.5281/zenodo.18986790}

\bibitem{collapse_emergence}
Lawson, E.R.\ (2026).
\textit{Mars Biosphere Collapse vs.\
Emergence --- Geometric Derivation
v1.0}.
Zenodo.
\href{https://doi.org/10.5281/zenodo.18995862}
{doi:10.5281/zenodo.18995862}

\bibitem{hurowitz2025}
Hurowitz, J.A., Tice, M.M., Allwood,
A.C., et al.\ (2025).
Redox-driven mineral and organic
associations in Jezero Crater, Mars.
\textit{Nature}, 645(8080), 332--340.
\href{https://doi.org/10.1038/s41586-025-09413-0}
{doi:10.1038/s41586-025-09413-0}

\bibitem{chivian2008}
Chivian, D., Brodie, E.L., Alm, E.J.,
et al.\ (2008).
Environmental genomics reveals a
single-species ecosystem deep within Earth.
\textit{Science}, 322(5899), 275--278.
\href{https://doi.org/10.1126/science.1155495}
{doi:10.1126/science.1155495}

\bibitem{orosei2018}
Orosei, R., Lauro, S.E., Pettinelli, E.,
et al.\ (2018).
Radar evidence of subglacial liquid
water on Mars.
\textit{Science}, 361(6401), 490--493.
\href{https://doi.org/10.1126/science.aar7268}
{doi:10.1126/science.aar7268}

\bibitem{insight2021}
Knapmeyer-Endrun, B., Panning, M.P.,
Bissig, F., et al.\ (2021).
Thickness and structure of the
Martian crust from InSight seismic data.
\textit{Science}, 373(6553), 438--443.
\href{https://doi.org/10.1126/science.abf8966}
{doi:10.1126/science.abf8966}

\bibitem{hirise2021}
Dundas, C.M., Bramson, A.M., Ojha, L.,
et al.\ (2021).
Exposed subsurface ice sheets in the
Martian mid-latitudes.
\textit{Science}, 359(6372), 199--201.
\href{https://doi.org/10.1126/science.aao1619}
{doi:10.1126/science.aao1619}

\bibitem{phoenix2009}
Smith, P.H., Tamppari, L.K., Arvidson,
R.E., et al.\ (2009).
H$_2$O at the Phoenix landing site.
\textit{Science}, 325(5936), 58--61.
\href{https://doi.org/10.1126/science.1172339}
{doi:10.1126/science.1172339}

\bibitem{odyssey2002}
Boynton, W.V., Feldman, W.C., Squyres,
S.W., et al.\ (2002).
Distribution of hydrogen in the near
surface of Mars: evidence for subsurface
ice deposits.
\textit{Science}, 297(5578), 81--85.
\href{https://doi.org/10.1126/science.1073722}
{doi:10.1126/science.1073722}

\bibitem{marsfield}
Acuna, M.H., Connerney, J.E.P.,
Ness, N.F., et al.\ (1999).
Global distribution of crustal
magnetization discovered by the Mars
Global Surveyor MAG/ER experiment.
\textit{Science}, 284(5415), 790--793.
\href{https://doi.org/10.1126/science.284.5415.790}
{doi:10.1126/science.284.5415.790}

\bibitem{ehlmann2010}
Ehlmann, B.L., Mustard, J.F.,
Murchie, S.L., et al.\ (2010).
Identification of hydrated silicate
minerals on Mars using MRO-CRISM.
\textit{Journal of Geophysical Research:
Planets}, 114(E2).
\href{https://doi.org/10.1029/2009JE003339}
{doi:10.1029/2009JE003339}

\bibitem{tharsis2023}
Hauber, E., Brož, P., Rossi, A.P.,
et al.\ (2023).
Long-lived and continual volcanic
eruptions and tectonic activity in
the Tharsis region of Mars.
\textit{Journal of Geophysical Research:
Planets}, 127(8).
\href{https://doi.org/10.1029/2022JE007511}
{doi:10.1029/2022JE007511}

\bibitem{mgc3catalog}
Cushing, G.E.\ (2015).
Candidate cave entrances on Mars.
\textit{Journal of Cave and Karst
Studies}, 74(1), 33--47.
\href{https://doi.org/10.4311/2010EX0167R}
{doi:10.4311/2010EX0167R}

\bibitem{webster2018}
Webster, C.R., Mahaffy, P.R.,
Atreya, S.K., et al.\ (2018).
Background levels of methane in Mars'
atmosphere show strong seasonal
variations.
\textit{Science}, 360(6393), 1093--1096.
\href{https://doi.org/10.1126/science.aaq0131}
{doi:10.1126/science.aaq0131}

\bibitem{moores2019}
Moores, J.E., King, P.L., Smith, C.L.,
et al.\ (2019).
The methane diurnal variation and
microseepage flux at Gale Crater, Mars
as constrained by the ExoMars Trace
Gas Orbiter and Curiosity observations.
\textit{Geophysical Research Letters},
46(16), 9430--9438.
\href{https://doi.org/10.1029/2019GL083800}
{doi:10.1029/2019GL083800}

\bibitem{magnetosomes}
Kopp, R.E.\ \& Kirschvink, J.L.\
(2008).
The identification and biogeochemical
interpretation of fossil magnetotactic
bacteria.
\textit{Earth-Science Reviews},
86(1--4), 42--61.
\href{https://doi.org/10.1016/j.earscirev.2007.08.001}
{doi:10.1016/j.earscirev.2007.08.001}

\bibitem{bibring2006}
Bibring, J.-P., Langevin, Y.,
Mustard, J.F., et al.\ (2006).
Global mineralogical and aqueous Mars
history derived from OMEGA/Mars Express
data.
\textit{Science}, 312(5772), 400--404.
\href{https://doi.org/10.1126/science.1122659}
{doi:10.1126/science.1122659}

\bibitem{boston2004}
Boston, P.J., Frederick, R.D.,
Welch, S.M., et al.\ (2004).
Extraterrestrial subsurface technology
test bed: human use and scientific
value of lava tubes on other worlds.
\textit{Advances in Space Research},
33(8), 1246--1252.
\href{https://doi.org/10.1016/S0273-1177(03)00538-5}
{doi:10.1016/S0273-1177(03)00538-5}

\bibitem{crater_darkening}
Daubar, I.J., McEwen, A.S., Byrne, S.,
Kennedy, M.R., \& Ivanov, B.\ (2013).
The current Martian cratering rate.
\textit{Icarus}, 225(1), 506--516.
\href{https://doi.org/10.1016/j.icarus.2013.04.009}
{doi:10.1016/j.icarus.2013.04.009}

\end{thebibliography}

% ── FOOTER ────────────────────────────────────────────────────
\vspace{1em}
\noindent\textcolor{rulecolor}{\rule{\linewidth}{0.4pt}}
\begin{center}
\footnotesize
\textit{The geometry came first.
The test comes after.
The record is here regardless of the outcome.}\\[0.3em]
Pre-registration chain:
\href{https://doi.org/10.5281/zenodo.18986790}
{10.5281/zenodo.18986790} (origin, 2026-03-12)
$\cdot$
Repository:
\href{https://github.com/Eric-Robert-Lawson/attractor-oncology}
{attractor-oncology}
$\cdot$
ORCID:
\href{https://orcid.org/0009-0002-0414-6544}
{0009-0002-0414-6544}
\end{center}

\end{document}
